
%-------------------------------------------------------------------------------
\documentclass[10pt]{article}
%-------------------------------------------------------------------------------

\usepackage{amscd}
\usepackage[T1]{fontenc}
\usepackage[spanish]{babel}
\usepackage{hyperref}
\usepackage[psamsfonts]{amssymb}
\usepackage[psamsfonts]{amssymb}
\usepackage{amsmath}
\usepackage[all]{xy}
\usepackage{graphicx}
\usepackage{float}
\usepackage{xspace}
\usepackage[utf8]{inputenc}
\usepackage{tikz}
\usetikzlibrary{fit,shapes.geometric}

%-------------------------------------------------------------------------------
\hoffset-1.0cm
\voffset-1.5cm
\textwidth6in
\textheight8.5in
%-------------------------------------------------------------------------------------------


%-------------------------------------------------------------------------------
\newcommand{\header}{\vspace*{-1cm} \hspace*{-1cm} \includegraphics[height=2cm, width=1.75cm]{comap.jpg} 
\centering 
	\begin{minipage}{12.5cm} \vspace{-2cm} \vskip4truept   \textbf{ {\Large VII CONGRESO NACIONAL DEL COLEGIO DE } }\\
	 \textbf{{\Large $\text{\ \ \ \ \ \ \ \ \ \ \ \ \  } $  MATEMATICOS DEL PERU}}\\
			\textbf{ {\large $\text{\   } $ Colegio de Matemáticos del Perú-Región Lambayeque} }\\
			 			 \centerline{\textrm{\small Del 1 al 3 y del 8 al 10 de octubre del 2025 }} 
	\end{minipage}
	\vskip1truept \hrule \vspace{0.2cm}
	\hspace*{0.0cm}}

\date{ }

%-------------------------------------------------------------------------------
\title{\header\\ \textbf{Schur multiplier of $\operatorname{SL}_2$ over finite commutative rings}}
%-------------------------------------------------------------------------------
\author{
	\textbf{Abraham Rojas Vega }%
	\thanks{e-mail: \texttt{abraham.rojas@usp.br}}%
	\\%
	ICMC-USP}
%-------------------------------------------------------------------------------
\begin{document}
	\maketitle
	%-------------------------------------------------------------------------------
	\begin{abstract}
		\noindent \textit{
		La segunda homología de $\operatorname{SL}_2$, también conocida como multiplicador de Schur, sobre anillos conmutativos finitos $A$ es el objeto de nuestra investigación. Probamos que $H_2 (\operatorname{SL}_2(A)) \cong K_2(A)$ siempre que el cuerpo residual $ A / \mathfrak{m}_A $ tenga característica impar diferente de $3, \ 5$ y $9$. Como aplicación, mostramos que si $\operatorname{GR}(p^l,m)$ es un anillo de Galois con cuerpo residuo $k$ de característica impar y $|k| \neq 3,5,9$, entonces $\operatorname{SL}_2(\operatorname{GR}(p^l,m))$ es trivial.
		}
		\vspace{2cm}
	\end{abstract}
	
	
	
	
	%Escriba las referencias más importantes para la obtención del trabajo
	\begin{thebibliography}{99}
		
		\bibitem{principal}
		Mirzaii, B., Rojas Vega, A.
		\textit{Schur multiplier of $\operatorname{SL}_2$ over finite commutative rings}.
		En preparación, 2025.

		\bibitem{torres,Mirzaii}
		Torres, A., Mirzaii, B.
		\textit{A refined scissors congruence group and the third homology of $\operatorname{SL}_2$}. Journal of Pure and Applied Algebra, \textbf{228}(6), 2024.
	\end{thebibliography}
	%-------------------------------------------------------------------------------
\end{document}
